\documentclass[aspectratio=169, table, xcdraw]{beamer}

\usetheme{Madrid}
\usecolortheme{default}
\definecolor{ucbblue}{RGB}{0, 51, 102}

\setbeamercolor{structure}{fg=ucbblue}
\setbeamercolor*{palette primary}{fg=white,bg=ucbblue}
\setbeamercolor*{palette secondary}{fg=white,bg=ucbblue!80}
\setbeamercolor*{palette tertiary}{fg=white,bg=ucbblue!90}
\setbeamercolor{title}{fg=white, bg=ucbblue}
\setbeamercolor{frametitle}{fg=white, bg=ucbblue}
\setbeamercolor{item}{fg=ucbblue}

\usepackage[T1]{fontenc}
\usepackage[utf8]{inputenc}
\usepackage{tgpagella}
\usefonttheme{serif}
\usepackage[spanish, es-noshorthands]{babel}
\usepackage{amsmath, amssymb, bm}
\usepackage{graphicx}
\usepackage[most]{tcolorbox}
\usepackage{booktabs}
\usepackage{tabularx}
\usepackage{tikz}
\usepackage{pgfplots}
\pgfplotsset{compat=1.18}
\usepackage[american]{circuitikz}
\usetikzlibrary{shapes.geometric, arrows.meta, positioning, calc}

\title[Pequeña Señal y Emisor Común]{Pequeña señal alrededor del punto Q y primer análisis de emisor común}
\subtitle{IMT-221 Circuitos Electrónicos III — Semana 6, Sesión 2}
\author[B. Quiroga Turdera]{M.Sc. Ing. Bernardo Quiroga Turdera}
\institute[UCB Tarija]{Universidad Católica Boliviana ``San Pablo'' — Sede Tarija}
\date{Semana 6}

\begin{document}

\begin{frame}
    \titlepage
\end{frame}

\begin{frame}{Principio Fundamental de la Sesión}
    \begin{tcolorbox}[colback=ucbblue!5, colframe=ucbblue, title=La Pequeña Señal NO reemplaza al Análisis DC]
        El modelo de pequeña señal \textbf{existe porque primero se estableció un punto de operación DC ($Q$)}. No tiene sentido construir un modelo incremental sin conocer el estado estático del transistor.
    \end{tcolorbox}
    
    \vspace{0.3cm}
    \textbf{Progresión de la clase:}
    \begin{center}
        $Q \rightarrow \text{DC + perturbación} \rightarrow \text{Linealización} \rightarrow g_m, r_e, r_\pi \rightarrow \text{Híbrido-}\pi \rightarrow \text{Equivalente AC} \rightarrow A_v$
    \end{center}
\end{frame}

\begin{frame}{Descomposición: DC + Pequeña Señal}
    Las señales totales instantáneas se componen del valor de polarización (DC) más una variación incremental (AC):
    \begin{align*}
        v_{BE}(t) &= V_{BEQ} + v_{be}(t) \\
        i_C(t) &= I_{CQ} + i_c(t) \\
        v_{CE}(t) &= V_{CEQ} + v_{ce}(t)
    \end{align*}
    
    \begin{center}
    \begin{tikzpicture}
        \begin{axis}[
            width=12cm, height=3.5cm,
            axis lines=middle,
            xmin=0, xmax=10, ymin=0, ymax=4,
            xtick=\empty, ytick={2}, yticklabels={$I_{CQ}$},
            xlabel={$t$}, ylabel={Corriente},
            every axis y label/.style={at={(ticklabel* cs:1.05)},anchor=south},
            every axis x label/.style={at={(ticklabel* cs:1.05)},anchor=west}
        ]
        % DC
        \addplot[thick, dashed, gray, domain=0:10] {2};
        % AC only (shifted for visualization)
        \addplot[thick, blue, domain=0:10, samples=100] {0.8*sin(deg(x*1.5)) + 0.5} node[pos=0.9, above] {$i_c(t)$};
        % Total
        \addplot[thick, red, domain=0:10, samples=100] {0.8*sin(deg(x*1.5)) + 2};
        \node[red] at (axis cs: 8, 3.2) {$i_C(t) = I_{CQ} + i_c(t)$};
        \end{axis}
    \end{tikzpicture}
    \end{center}
    \textit{Pequeña señal no significa baja frecuencia; significa perturbación pequeña alrededor de $Q$.}
\end{frame}

\begin{frame}{Linealización Local en el Punto Q}
    \begin{columns}
        \begin{column}{0.48\textwidth}
            La relación fundamental es exponencial:
            \begin{equation*}
                i_C = I_S e^{v_{BE}/V_T}
            \end{equation*}
            Si la excursión de $v_{be}$ es muy pequeña ($|v_{be}| \ll V_T$), la curva puede aproximarse por su recta tangente en $Q$.
            
            \vspace{0.2cm}
            La \textbf{transconductancia} $g_m$ es esa pendiente:
            \begin{equation*}
                g_m = \left. \frac{\partial i_C}{\partial v_{BE}} \right|_Q = \frac{I_{CQ}}{V_T}
            \end{equation*}
        \end{column}
        \begin{column}{0.48\textwidth}
            \centering
            \begin{tikzpicture}
                \begin{axis}[
                    width=7cm, height=5.5cm,
                    axis lines=left,
                    xmin=0.5, xmax=0.8, ymin=0, ymax=5,
                    xlabel={$v_{BE}$ (V)}, ylabel={$i_C$ (mA)},
                    xtick={0.7}, xticklabels={$V_{BEQ}$},
                    ytick={2}, yticklabels={$I_{CQ}$},
                    grid=none
                ]
                % Exponencial simulada
                \addplot[thick, blue, domain=0.5:0.75, samples=50] {0.001 * exp(10.85*x)};
                % Punto Q
                \addplot[mark=*, red] coordinates {(0.7, 1.99)} node[above left] {$Q$};
                % Tangente
                \addplot[thick, red, dashed, domain=0.65:0.75] {21.6*(x-0.7) + 1.99};
                \node at (axis cs: 0.74, 1) {$g_m = \text{pendiente}$};
                \end{axis}
            \end{tikzpicture}
        \end{column}
    \end{columns}
\end{frame}

\begin{frame}{Parámetros Incrementales del BJT}
    Asumiendo $T \approx 300\,\text{K} \implies V_T \approx 25.85\,\text{mV}$. El modelo incremental de 1\textsuperscript{er} orden define:
    
    \vspace{0.2cm}
    \begin{block}{1. Transconductancia ($g_m$)}
        \begin{equation*}
            g_m = \frac{I_{CQ}}{V_T} \quad [\text{S}]
        \end{equation*}
        Controla la fuente de corriente dependiente ($i_c = g_m v_{be}$).
    \end{block}
    
    \begin{columns}
        \begin{column}{0.48\textwidth}
            \begin{block}{2. Resistencia intrínseca ($r_e$)}
                \begin{equation*}
                    r_e \approx \frac{V_T}{I_{CQ}} \approx \frac{1}{g_m} \quad [\Omega]
                \end{equation*}
                \textit{No es un resistor físico externo.}
            \end{block}
        \end{column}
        \begin{column}{0.48\textwidth}
            \begin{block}{3. Resistencia de Base ($r_\pi$)}
                \begin{equation*}
                    r_\pi = \frac{\beta}{g_m} \quad [\Omega]
                \end{equation*}
                Resistencia vista desde la base al emisor.
            \end{block}
        \end{column}
    \end{columns}
\end{frame}

\begin{frame}{Construcción del Modelo Híbrido-$\pi$}
    Reemplazamos el BJT no lineal por una red lineal válida \textit{sólo} para pequeñas variaciones AC. \\
    \textbf{Suposición de enseñanza:} Efecto Early despreciado ($r_o \rightarrow \infty$).
    
    \vspace{0.4cm}
    \begin{columns}
        \begin{column}{0.3\textwidth}
            \centering
            \textbf{Símbolo NPN}
            \vspace{0.2cm}
            
            \begin{circuitikz}[scale=1.2, transform shape, thick]
                \draw (0,0) node[npn](Q){};
                \node[above] at (Q.C) {C};
                \node[below] at (Q.E) {E};
                \node[left]  at (Q.B) {B};
                \draw[->, red] (-0.5,-0.1) arc (-180:-90:0.4) node[midway, left]{$v_{be}$};
            \end{circuitikz}
        \end{column}
        \begin{column}{0.1\textwidth}
            \centering
            $\implies$
        \end{column}
        \begin{column}{0.5\textwidth}
            \centering
            \textbf{Modelo Híbrido-$\pi$}
            \vspace{0.2cm}
            
            \begin{circuitikz}[scale=1, transform shape, thick]
                \draw (0,0) node[left]{B} to[R, l=$r_\pi$, v=$v_\pi$] (0,-2.5) node[below]{E};
                \draw (3,0) node[right]{C} to[cI, l=$g_m v_\pi$] (3,-2.5) -- (0,-2.5);
                \draw (0,0) -- (0.5,0);
                \draw (3,0) -- (2.5,0);
            \end{circuitikz}
        \end{column}
    \end{columns}
\end{frame}

\begin{frame}{Regla de la Tierra AC (AC Ground)}
    \begin{alertblock}{El Principio de Superposición en AC}
        Una fuente de tensión DC ideal tiene una variación incremental de cero ($\Delta V_{CC} = 0$). En el análisis de pequeña señal, sus terminales se conectan a \textbf{Tierra AC}.
    \end{alertblock}
    
    \begin{columns}
        \begin{column}{0.45\textwidth}
            \centering
            \textbf{Circuito Original DC}
            \begin{circuitikz}[scale=0.8, transform shape, thick]
                \draw (0,3) -- (3,3) node[above]{$+V_{CC}$};
                \draw (0,3) to[R, l_={$R_B$}] (0,1) node[below]{B};
                \draw (3,3) to[R, l={$R_C$}] (3,1) node[below]{C};
            \end{circuitikz}
        \end{column}
        \begin{column}{0.1\textwidth}
            $\implies$
        \end{column}
        \begin{column}{0.45\textwidth}
            \centering
            \textbf{Red Externa AC}
            \begin{circuitikz}[scale=0.8, transform shape, thick]
                \draw (0,3) node[ground, rotate=180]{} -- (0,2);
                \draw (3,3) node[ground, rotate=180]{} -- (3,2);
                \draw (0,3) to[R, l_={$R_B$}] (0,0) node[below]{B};
                \draw (3,3) to[R, l={$R_C$}] (3,0) node[below]{C};
            \end{circuitikz}
        \end{column}
    \end{columns}
    \vspace{0.2cm}
    \textit{Consecuencia vital: ¡Los resistores de polarización NO desaparecen! Cambian de nodo de referencia.}
\end{frame}

\begin{frame}{Transformación del Emisor Común (Paso a Paso)}
    \centering
    \begin{columns}
        \begin{column}{0.33\textwidth}
            \centering
            \textbf{1. Original}
            \resizebox{\textwidth}{!}{
            \begin{circuitikz}[thick]
                \draw (0,0) node[npn](Q){};
                \draw (Q.E) node[ground]{};
                \draw (Q.C) to[R, l_={$R_C$}] (0,2) -- (-2,2) node[above]{$+V_{CC}$};
                \draw (Q.B) -- (-2,0) to[R, l={$R_B$}] (-2,2);
                \draw (-2.5,0) node[left]{$v_i$} to[short, o-] (-2,0);
                \draw (Q.C) to[short, -o] (1,0.77) node[right]{$v_o$};
            \end{circuitikz}
            }
        \end{column}
        \begin{column}{0.33\textwidth}
            \centering
            \textbf{2. Externa AC}
            \resizebox{\textwidth}{!}{
            \begin{circuitikz}[thick]
                \draw (0,0) node[npn](Q){};
                \draw (Q.E) node[ground]{};
                \draw (Q.C) to[R, l_={$R_C$}] (0,2) node[ground, rotate=180]{};
                \draw (Q.B) -- (-2,0) to[R, l={$R_B$}] (-2,2) node[ground, rotate=180]{};
                \draw (-2.5,0) node[left]{$v_i$} to[short, o-] (-2,0);
                \draw (Q.C) to[short, -o] (1,0.77) node[right]{$v_o$};
            \end{circuitikz}
            }
        \end{column}
        \begin{column}{0.33\textwidth}
            \centering
            \textbf{3. Híbrido-$\pi$}
            \resizebox{\textwidth}{!}{
            \begin{circuitikz}[thick]
                \draw (0,0) to[R, l_={$R_B$}] (0,2) node[ground, rotate=180]{};
                \draw (1,0) to[R, l_={$r_\pi$}, v^={$v_\pi$}] (1,-2) node[ground]{};
                \draw (-1,0) node[left]{$v_i$} to[short, o-] (1,0);
                
                \draw (3,0) to[cI, l={$g_m v_\pi$}] (3,-2) node[ground]{};
                \draw (4,0) to[R, l={$R_C$}] (4,2) node[ground, rotate=180]{};
                \draw (3,0) -- (4,0);
                \draw (4,0) to[short, -o] (5,0) node[right]{$v_o$};
            \end{circuitikz}
            }
        \end{column}
    \end{columns}
\end{frame}

\begin{frame}{Derivación de la Ganancia ($A_v$) y Fase}
    Del modelo Híbrido-$\pi$ previo, observamos la entrada y la salida:
    \begin{enumerate}
        \item Como el emisor está a tierra AC: $\quad \mathbf{v_\pi = v_i}$
        \item La fuente dependiente inyecta corriente hacia tierra, extrayéndola del nodo colector:
        \begin{equation*}
            i_c = g_m v_\pi
        \end{equation*}
        \item Esta corriente fluye por $R_C$ desde tierra hacia el colector (KCL), induciendo una caída negativa:
        \begin{equation*}
            v_o = -i_c R_C = -g_m v_\pi R_C
        \end{equation*}
    \end{enumerate}
    
    \begin{block}{Ganancia de Tensión Idealizada ($r_o \to \infty$, sin carga externa)}
        \begin{equation*}
            A_v = \frac{v_o}{v_i} = -g_m R_C
        \end{equation*}
    \end{block}
    \textit{El signo negativo significa inversión de fase (180$^\circ$), no atenuación.}
\end{frame}

\begin{frame}{Efecto de Carga y Límite de Excursión Lineal}
    \begin{columns}
        \begin{column}{0.5\textwidth}
            \textbf{Efecto de la Carga ($R_L$):} \\
            Si acoplamos una resistencia de carga en AC al colector, ésta queda en paralelo con $R_C$:
            \begin{equation*}
                A_v \approx -g_m (R_C \parallel R_L)
            \end{equation*}
            La ganancia magnitud disminuye invariablemente al agregar carga.
        \end{column}
        \begin{column}{0.48\textwidth}
            \textbf{Precaución de Rango Lineal:} \\
            Un gran $|A_v|$ teórico no garantiza amplificación ilimitada. La excursión total $V_{CE}(t)$ está atrapada entre el \textbf{Corte} y la \textbf{Saturación}.
            \begin{center}
            \begin{tikzpicture}[scale=0.7]
                \draw[thick, ->] (0,0) -- (5,0) node[right]{$t$};
                \draw[thick, ->] (0,-1) -- (0,3) node[above]{$v_{CE}(t)$};
                \draw[dashed, red] (0,2.5) -- (5,2.5) node[right]{Corte};
                \draw[dashed, red] (0,0.5) -- (5,0.5) node[right]{Sat.};
                \draw[blue, thick, domain=0:4.5, samples=100] plot (\x, {1.5 + 1.2*sin(deg(\x*3))});
                \node at (2.5, 3) {\small ¡Recorte si la señal es muy grande!};
            \end{tikzpicture}
            \end{center}
        \end{column}
    \end{columns}
\end{frame}

\begin{frame}{Puente hacia el Hito 2: El Par Diferencial}
    ¿Hacia dónde nos dirige este análisis de un solo transistor?
    
    \vspace{0.4cm}
    \centering
    \begin{tikzpicture}[node distance=2.5cm, auto, thick]
        \node[draw, rounded corners, fill=blue!10] (bjt) {Un solo BJT};
        \node[draw, rounded corners, fill=blue!10, right of=bjt, xshift=1cm] (ce) {Emisor Común};
        \node[draw, rounded corners, fill=red!10, right of=ce, xshift=1.5cm] (diff) {\textbf{Par Diferencial}};
        
        \draw[->] (bjt) -- node[above]{$g_m, r_\pi$} (ce);
        \draw[->] (ce) -- node[above]{Acoplar 2 BJTs} (diff);
    \end{tikzpicture}
    
    \vspace{0.5cm}
    \begin{itemize}
        \item El amplificador diferencial usa la misma física incremental que hemos desarrollado hoy.
        \item Entender $g_m$ y la referencia de emisor es el prerrequisito directo para analizar el modo diferencial y el rechazo en modo común (CMRR).
    \end{itemize}
\end{frame}

\end{document}
